\documentclass[11pt]{article}

\usepackage[margin=1in]{geometry}
\usepackage{amsmath,amssymb,amsthm}
\usepackage{mathtools}
\usepackage{enumitem}
\usepackage{hyperref}
\usepackage{cite}

\hypersetup{
  colorlinks=true,
  linkcolor=black,
  citecolor=black,
  urlcolor=blue
}

\newtheorem{definition}{Definition}
\newtheorem{remark}{Remark}
\newtheorem{proposition}{Proposition}

\title{Deriving an Implication Frame from an LLM}
\author{Bradley P.\ Allen\\University of Amsterdam}
\date{Technical note, April 2026}

\begin{document}

\maketitle

\begin{abstract}
Hlobil and Brandom's implication-space semantics takes the material consequence relation $I$ over a set of bearers $B$ as given. This note shows how to derive, from a large language model $M$, an implication frame $\langle B, I_M \rangle$ via a selection function applied to $M$'s bearer expression probabilities. The frame satisfies Containment by construction and leaves the framework's remaining structural commitments open.
\end{abstract}

\section{Setup}

Let $V$ be a finite set of tokens and $\mathcal{L} = V^*$ the set of finite token sequences. A non-empty $B \subseteq \mathcal{L}$ is the set of \emph{bearers}.

The construction below works within the implication-space semantics of Hlobil and Brandom \cite{hlobil2025}, to which the reader is referred for the following machinery used without restatement: the implication frame $\langle B, I \rangle$ with $I \subseteq \wp(B) \times \wp(B)$, $\langle \emptyset, \emptyset \rangle \notin I$, and $\langle B, B \rangle \in I$ \cite[Def.~62]{hlobil2025}; the range of subjunctive robustness $\mathrm{RSR}$ \cite[Def.~63]{hlobil2025}; implication equivalence $\approx$ \cite[Def.~64]{hlobil2025}; implicational role $R(\alpha) = \langle R^+(\alpha), R^-(\alpha) \rangle$ \cite[Def.~65]{hlobil2025}; conceptual content as a pair of roles \cite[Def.~66]{hlobil2025}; the NMMS clauses for the connectives \cite[Def.~70]{hlobil2025}; and Corollary~77, which establishes that a Containment-satisfying frame extended via NMMS yields classical propositional logic in the narrowly logical fragment of the extended consequence relation.

An implication frame $\langle B, I \rangle$ \emph{obeys Containment} iff $\Gamma \cap \Delta \neq \emptyset$ implies $\langle \Gamma, \Delta \rangle \in I$. Containment is the only structural constraint the construction below will impose.

\section{The Derivation}

\begin{definition}[Bearer expression probabilities of an LLM]
\label{def:mu}
Let $M$ be a large language model with completion distribution $p_M : \mathcal{L} \to \Delta(\mathcal{L})$ assigning to each context $c \in \mathcal{L}$ a probability distribution over token sequences, and let $\mathrm{Expr} \subseteq \mathcal{L} \times B$ be an expression relation such that $\langle \sigma, \varphi \rangle \in \mathrm{Expr}$ iff $\sigma$ expresses $\varphi$. For each $\varphi \in B$, let
\[
E_\varphi = \{\sigma \in \mathcal{L} : (\sigma, \varphi) \in \mathrm{Expr}\}
\]
be the event that a completion expresses $\varphi$. The \emph{bearer expression probabilities of $M$} are given by the function $\mu_M : \mathcal{L} \times B \to [0,1]$ where
\[
\mu_M(c)(\varphi) \;=\; p_M(c)(E_\varphi) \;=\; \sum_{\sigma \in E_\varphi} p_M(c)(\sigma).
\]
That is, $\mu_M(c)(\varphi)$ is the probability under $p_M(c)$ that $M$'s completion of $c$ expresses $\varphi$.
\end{definition}

\begin{remark}[Expression events]
\label{rem:expression-events}
Each $\mu_M(c)(\varphi)$ is the probability of a separate event in $\mathcal{L}$ under $p_M(c)$. The events $E_\varphi$ may overlap (a single $\sigma$ may express several bearers) and need not cover $\mathcal{L}$ (some $\sigma$ may express no bearer in $B$), so $\sum_{\varphi \in B} \mu_M(c)(\varphi)$ may exceed or fall short of $1$.
\end{remark}

For $\Gamma \subseteq B$, let $\mathrm{Ctx}_\Gamma = \{c \in \mathcal{L} : (c, \gamma) \in \mathrm{Expr} \text{ for all } \gamma \in \Gamma\}$ denote the set of contexts that express every bearer in $\Gamma$. A \emph{selection function} is a map $S$ that takes a function $f : B \to [0,1]$ to a subset $S(f) \subseteq B$. Fix throughout an admissible selection function $S$; we write $I_M$ for the frame so derived.

\begin{definition}[Derived implication frame]
\label{def:frame}
Given an LLM $M$ with bearer expression probabilities $\mu_M$ and a selection function $S$, the \emph{derived implication frame for $M$} is $\langle B, I_M \rangle$ where $\langle \Gamma, \Delta \rangle \in I_M$ iff at least one of the following holds:
\begin{enumerate}[label=(\roman*)]
\item $\Gamma \cap \Delta \neq \emptyset$; or
\item $\exists c \in \mathrm{Ctx}_\Gamma$ such that $\Delta \cap S(\mu_M(c)) \neq \emptyset$.
\end{enumerate}
\end{definition}

\begin{definition}[Admissible selection functions]
\label{def:selection}
The construction admits a range of selection functions; three natural instances are:
\begin{itemize}
\item \emph{Top-$k$}: for $k \geq 1$, $\mathrm{Top}_k(\mu_M(c))$ is the set of bearers $\varphi$ with $\mu_M(c)(\varphi) > 0$ among the $k$ with largest $\mu_M(c)(\varphi)$, ties broken by a fixed convention on $B$.
\item \emph{Absolute threshold}: for $\tau \in (0, 1]$, $\mathrm{Thr}_\tau(\mu_M(c)) = \{\varphi \in B : \mu_M(c)(\varphi) > \tau\}$.
\item \emph{Nucleus} (top-$p$): for $p \in (0, 1]$, ordering bearers $\varphi_1, \varphi_2, \ldots$ by descending $\mu_M(c)(\varphi)$ with ties broken by a fixed convention, $\mathrm{Nuc}_p(\mu_M(c)) = \{\varphi_1, \ldots, \varphi_n\}$ where $n$ is the smallest index with $\sum_{i=1}^n \mu_M(c)(\varphi_i) \geq p$, or all bearers with positive mass if no such $n$ exists. Cumulative mass is taken in absolute terms; by Remark~\ref{rem:expression-events}, $\mu_M(c)$ is not a pmf on $B$.
\end{itemize}
\end{definition}

Clause (ii) of Definition~\ref{def:frame} is non-trivial under any of these instances: it requires positive support from $M$ for some $\delta \in \Delta$, strictly stronger than ``the completion is possible'' (vacuous under positive-temperature sampling) and, when $S(\mu_M(c))$ is non-singleton, strictly weaker than ``the completion is the single most probable output''.

\begin{remark}[Methodological commitment of $S$]
\label{rem:methodological}
The choice of selection function is a substantive methodological commitment the construction itself is silent on. Top-$k$ is engineering convenience: $k$ is chosen for tractability and has no independent principled basis. Absolute threshold operationalizes the Lockean thesis from formal epistemology---that credence above a fixed threshold counts as commitment \cite{foley1992}---and inherits its known difficulties (the lottery and preface paradoxes). Nucleus matches the regime under which deployed LLMs generate text \cite{holtzman2020}, but adapts cardinality to context rather than answering to an independent epistemic norm. The downstream structural properties of $I_M$, and the content structure derived from it, depend on $S$.
\end{remark}

\begin{remark}[Stratification]
\label{rem:stratification}
Definition~\ref{def:frame} does not refer to $R$, $\approx$, or content. The content structure that arises from applying the Hlobil--Brandom construction to $I_M$ is therefore a consequence of the derivation, not a presupposition. Circularity that would arise from defining $I_M$ via content-equivalence of contexts is blocked at the bearer level.
\end{remark}

\begin{remark}[Containment by construction]
\label{rem:containment}
By clause (i) of Definition~\ref{def:frame}, $\langle B, I_M \rangle$ obeys Containment.
\end{remark}

\section{Discussion}

With $\langle B, I_M \rangle$ in hand, the standard Hlobil--Brandom machinery applies without modification: $\mathrm{RSR}$, $\approx$, implicational roles, conceptual contents, and extensions of the material base with logical vocabulary are all obtained by their constructions operating on $\langle B, I_M \rangle$ as an ordinary implication frame. Nothing specific to the LLM origin of $I_M$ enters these constructions.

\begin{proposition}[Classical core under NMMS]
\label{prop:classical}
Let $\langle B, I_M \rangle$ be a derived implication frame. The narrowly logical fragment of the consequence relation obtained by extending $\langle B, I_M \rangle$ via the NMMS clauses is classical propositional logic.
\end{proposition}

\begin{proof}
By Remark~\ref{rem:containment}, $\langle B, I_M \rangle$ satisfies Containment. The result follows from Corollary~77 of Hlobil and Brandom \cite{hlobil2025}.
\end{proof}

The classical core holds for any $M$ and any admissible parameters; it is a property of the NMMS-extended consequence relation, not of $I_M$. The material base may be substructural while the extension carries a classical logical core: whether $I_M$ is in fact nonmonotonic, nontransitive, or non-contractive is an empirical property of $\mu_M$, not something the construction forces or forbids.

Implicational role inclusion and content inclusion \cite[Def.~90, 93]{hlobil2025} are likewise available in the derived frame: one can ask which bearers $M$ treats as intersubstitutable \emph{salva consequentia} and which contents are included in which others. These questions have determinate answers given $\langle B, I_M \rangle$ and characterize aspects of $M$'s behavior not directly visible in $\mu_M$.

The construction thus provides a determinate answer to the question of LLM content: given $M$, a bearer set $B$, and admissible parameters, the content structure derived from $\langle B, I_M \rangle$ is well-defined and characterizes $M$'s material-inferential dispositions at the level of role and content identity.

\begin{thebibliography}{9}

\bibitem{hlobil2025}
Hlobil, U., \& Brandom, R.~B. (2025). \emph{Reasons for Logic, Logic for Reasons: Pragmatics, Semantics, and Conceptual Roles}. Routledge.

\bibitem{foley1992}
Foley, R. (1992). The epistemology of belief and the epistemology of degrees of belief. \emph{American Philosophical Quarterly}, 29(2), 111--124.

\bibitem{holtzman2020}
Holtzman, A., Buys, J., Du, L., Forbes, M., \& Choi, Y. (2020). The curious case of neural text degeneration. In \emph{International Conference on Learning Representations}.

\end{thebibliography}

\end{document}
